% Context from chapters-tex/approximation.tex; for mathematical reference, not compilation.
f scales allow us to define a comparison approximation
\eql{\label{eq-hand-made-approx}
	\tilde f_M=P_{V_0}f+\sum_{j_2\leq j\leq0} \sum_{n \in \Cc_j} \dotp{f}{\psi_{j,n}} \psi_{j,n} +  \sum_{j_1\leq j\leq0} \sum_{n \in \Cc_j^c} \dotp{f}{\psi_{j,n}} \psi_{j,n}.
}
The error of the comparison $M$-term approximation $\tilde f_M$ bounds the best $M$-term error from above:
\eqml{\label{eq-wavapprox-1d-1}
	\norm{f-f_M}^2 &\leq \norm{f-\tilde f_M}^2 \leq \sum_{j<j_2, n \in \Cc_j} |\dotp{f}{\psi_{j,n}}|^2 + 
\sum_{j<j_1, n \in \Cc_j^c} |\dotp{f}{\psi_{j,n}}|^2 \\
	& \leq \sum_{j < j_2} (K |\Ss|) \times C^2 2^{j} + \sum_{j < j_1} 2^{-j} \times  C^2 2^{j(2\al+1)} \\
	& = O( 2^{j_2} + 2^{2\al j_1} ) = O( T^2 + T^{ \frac{2\al}{\al+1/2} } ) = O( T^{ \frac{2\al}{\al+1/2} } ).
}

  
\paragraph{Step 3. Counting the Coefficients.}
The number of coefficients needed to build the approximating signal $\tilde f_M$ is
\eqml{\label{eq-wavapprox-1d-2}
	M &\leq1+\sum_{j_2\leq j\leq0} |\Cc_j| +
\sum_{j_1\leq j\leq0} |\Cc_j^c| \leq
1+\sum_{j_2\leq j\leq0} K |\Ss| +
\sum_{j_1\leq j\leq0} 2^{-j} \\
	&= O( |\log(T)| + T^{\frac{-1}{\al+1/2}} ) 
=  O(T^{\frac{-1}{\al+1/2}} ).
}
  
\paragraph{Step 4. Putting everything together.}
Given a coefficient budget $M$, choose $T$ to be a sufficiently large constant times $M^{-(\alpha+1/2)}$. The coefficient-count bound then uses at most $M$ terms, and the error bound is $O(M^{-2\alpha})$.
\end{proof}

For piecewise Lipschitz signals, the wavelet rate is faster than the jump-limited $O(M^{-1})$ Fourier bound in~\eqref{eq-fourier-piecewisesmooth}. It also matches the smooth-signal rate~\eqref{eq-sovolev-signal-decay}: finitely many singularities do not worsen the asymptotic exponent in one dimension. The rate~\eqref{eq-decay-piecewav} is asymptotically optimal for this class.

\myfigure{
\image{approximation}{.6}{wav-approx-1d-1}\\
\image{approximation}{.6}{wav-approx-1d-2}\\
\image{approximation}{.6}{wav-approx-1d-3}
}{ 
	1-D wavelet approximation.  
}{fig-wav-approx-1d}

Figure \ref{fig-wav-approx-1d} shows examples of wavelet approximation of singular signals. 


                                        
\subsection{2-D Piecewise Smooth Approximation}

We now give the counterpart of Theorem~\ref{thm-approx-piecesmooth-1d} for 2-D functions.

\begin{thm}\label{thm-approx-piecesmooth-2d}
Let $f$ belong to the piecewise smooth image model~\eqref{eq-model-piece-smooth-2d}, with $\alpha\geq1$ and finitely many rectifiable edge curves. For compactly supported orthonormal wavelets with at least $\alpha$ vanishing moments and compatible boundary treatment, the best $M$-term error satisfies
\eql{\label{eq-approximation-error-wav-2d}
	\norm{f-f_M}^2 = O(M^{-1}).
}
\end{thm}

\begin{proof}
For an image in the model~\eqref{eq-model-piece-smooth-2d}, define the singular index set $\Cc_j$ as in~\eqref{eq-singsupport-1d}. Unlike the 1-D case, its cardinality grows as the scale $2^j$ decreases:
\eq{
	|\Cc_j^\om| \leq 2^{-j} K |\Ss|,
}
where $|\Ss|$ is the total length of the singular curves $\Ss$, and $\om \in \{V,H,D\}$ is the wavelet orientation.
Using \eqref{eq-wavcoef-decay-1} for $d=2$, the decay of regular coefficients is bounded as
\eq{
	\foralls n \in (\Cc_j^{\om})^c, \quad \abs{\dotp{f}{\psi_{j,n}^\om}} \leq  C2^{j(\alpha+1)}.
}
Using \eqref{eq-wavcoef-decay-2} for $d=2$, the decay of singular coefficients is bounded as
\eq{
	\foralls n \in \Cc_j^\om, \quad \abs{\dotp{f}{\psi_{j,n}^\om}} \leq  C 2^{j}.
}
After fixing $T$, the cut-off scales are defined as 
\eq{
	2^{j_1} = (T/C)^{\frac{1}{\al+1}}
	\qandq
	2^{j_2} = T/C.
}
We define, as in \eqref{eq-hand-made-approx}, a comparison approximation.
Similarly to \eqref{eq-wavapprox-1d-1}, we bound the approximation error as
\eq{
	\norm{f-f_M}^2 \leq \norm{f-\tilde f_M}^2 
	= O( 2^{j_2} + 2^{2\al j_1} ) = O( T + T^{ \frac{2\al}{\al+1} } ) = O( T )
}
and the number of coefficients as 
\eq{
	M=O(T^{-1}+T^{-2/(\alpha+1)}) = O(T^{-1} ).
}
Since $\alpha\geq1$, choose $T$ of order $M^{-1}$ with a sufficiently large constant. This meets the coefficient budget and gives the asserted error. For $0<\alpha<1$, the same argument gives $O(M^{-\alpha})$ instead.
\end{proof}

The wavelet rate improves on the $O(M^{-1/2})$ Fourier rate~\eqref{eq-fourier-decay-2d}, but it still fails to exploit the full $\Cal$ regularity away from the edge curves.

This rate also holds for the broader bounded-variation model~\eqref{eq-model-tv}, for which wavelet approximation is asymptotically optimal.

\myfigure{
\tabtrois{
\image{approximation}{.32}{approx-2d-wav-1}&
\image{approximation}{.32}{approx-2d-wav-2}&
\image{approximation}{.32}{approx-2d-wav-3}
}
}{ 
	2-D wavelet approximation.  
}{fig-approx-2d-wav}

Figure \ref{fig-approx-2d-wav} shows wavelet approximations of a bounded variation image. 

                                                                                                                                                                 


                                        
                                                                                                                                                                                                                             

                                           
                                                                                                                                                                                               

                                        
                                        
\section{Approximation of Cartoon Images}

The isotropic, compactly supported wavelets considered here occupy squares that do not follow the smooth contours of geometric images~\eqref{eq-cartoon-model}. These contours have more structure than the finite-length level sets allowed by the bounded-variation model~\eqref{eq-model-tv}.

                                        
\subsection{Wavelet Approximation of Cartoon Images}

The bound~\eqref{eq-approximation-error-wav-2d} gives a wavelet approximation rate of $O(M^{-1})$ for cartoon images~\eqref{eq-cartoon-model}. Even the simple indicator $f=1_\Om$, where $\Om$ is a disk, can attain this slow rate: isotropic wavelets do not efficiently follow its curved boundary. The smoothed cartoon model~\eqref{eq-cartoon-model-smooth} can behave similarly when the coefficient budget $M$ is too small to resolve the blur.

Figure \ref{fig-cartoon-wav} shows that many large coefficients are located near edge curves, and retaining only a small number leads to a poor approximation with visually unpleasant artifacts.


\myfigure{
\tabtrois{
\image{approximation}{.32}{cartoon-wav-image}&
\image{approximation}{.